\begin{frame}[t]{\ev{\tagO}Checkpoint success is conditional on reaching the checkpoint}
\begin{center}
\begin{tikzpicture}[x=1cm,y=1cm]
\node[state,minimum width=2.1cm] (a) at (0,0) {task start};
\node[state,minimum width=2.6cm] (b) at (4.6,0) {reached checkpoint};
\node[candidate,minimum width=2.6cm,minimum height=.6cm] (c) at (9.6,.7) {continuation A};
\node[candidate,minimum width=2.6cm,minimum height=.6cm] (d) at (9.6,0) {continuation B};
\node[candidate,minimum width=2.6cm,minimum height=.6cm] (e) at (9.6,-.7) {continuation C};
\draw[flow] (a)--node[above,font=\scriptsize]{expensive prefix}node[below,font=\scriptsize,text=Teal]{reaches it: 50\%}(b);
\draw[flow] (b)--(c);\draw[flow] (b)--(d);\draw[flow] (b)--(e);
\node[font=\scriptsize,text=Teal] at (9.6,1.25) {each succeeds given it: 80\%};
\end{tikzpicture}
\end{center}
{\small\textbf{Illustration.}\quad end-to-end success $=0.5\times0.8=0.4$, while a checkpoint-started run shows 0.8.\par\vspace{1pt}
\begin{ticks}
\item A checkpoint distribution differs from the task-start distribution.
\item When every repair inherits one plan, successful continuations cannot validate that plan.
\item Go-Explore makes reuse of reached states explicit: first return, then explore.
\end{ticks}\par}
\claim{Successes after a checkpoint say nothing about the prefix they share.}
\sources{\href{https://www.nature.com/articles/s41586-020-03157-9}{Ecoffet et al., First return, then explore, Nature (2021)}. Probabilities are illustrative, not measured.}
\note{[0:35] Success after a checkpoint is conditional on reaching it. In this illustration, the prefix reaches a usable state half the time. A continuation succeeds eighty percent of the time. End-to-end success is forty percent. Go-Explore makes this reuse of reached states explicit. When every repair inherits one plan, their successes cannot validate that plan. Next, we define what a checkpoint captured.}
\end{frame}
